\chapter{Минимальная архитектурная развилка после запрета \texorpdfstring{$R_2$}{R2}}
\label{ch:v7-r2-minimal-architecture-branch-gate}

\section{Вопрос}

Строгий первый порядок запретил прямой $R_2$ на фиксированном $H_{15}$.
Теперь необходимо сравнить три минимальных изменения архитектуры:
автоморфную скрутку неизменённой алгебры, обобщённые внутренние флуктуации
без первого порядка и строгий путь с новыми фермионными вершинами.

\section{Автоморфная скрутка не меняет типы слагаемых}

Стандартная алгебра разлагается на попарно неизоморфные простые вещественные
идеалы
\begin{equation}
 \mathcal A_{\rm SM}=\mathbb C\oplus\mathbb H\oplus M_3(\mathbb C).
 \label{eq:v7-r2-branch-simple-ideals}
\end{equation}
Любой $*$-автоморфизм прямой суммы может переставлять только изоморфные
простые идеалы. Поэтому
\begin{equation}
 \rho(\mathbb C)=\mathbb C,\qquad
 \rho(\mathbb H)=\mathbb H,\qquad
 \rho(M_3)=M_3.
 \label{eq:v7-r2-branch-automorphism-orbits}
\end{equation}
Внутренние автоморфизмы могут вращать матричные координаты внутри $\mathbb H$
и $M_3$, но не превращают метку $\mathbb C$ в $M_3$ или $\mathbb H$.
Следовательно, простая автоморфная скрутка неизменённой $\mathcal A_{\rm SM}$
не превращает диагональ $R_2$ в ребро одной строки или одного столбца.
Общие теории скрученных Real-троек существуют, но применение скрутки требует
явного расширения представленной алгебры или другой реализации
\cite{v7r2branch-landi,v7r2branch-grand}.

\section{Нулевой по новым вершинам путь: отказ от первого порядка}

Для Real-тройки без условия первого порядка к линейным флуктуациям добавляется
нелинейный член \cite{v7r2branch-generalized}:
\begin{equation}
 D_A=D+A_{(1)}+\widetilde A_{(1)}+A_{(2)},
 \label{eq:v7-r2-branch-generalized-dirac}
\end{equation}
где
\begin{equation}
 A_{(2)}=\sum_j\widehat a_j[A_{(1)},\widehat b_j]
 =\sum_{j,k}\widehat a_j a_k[[D,b_k],\widehat b_j].
 \label{eq:v7-r2-branch-quadratic-fluctuation}
\end{equation}
При строгом первом порядке $A_{(2)}=0$; при его нарушении именно этот член
восстанавливает нелинейную калибровочную ковариантность полной флуктуации.

Этот маршрут не требует новых фермионных вершин, но меняет аксиоматику и
действие. Из общей формулы ещё не следует, что на конкретном $H_{15}$
возникает ровно $R_2$, что его нормировка единственна или что вакуум
сохраняет цвет.

\section{Строгий путь: минимальный зеркальный цикл}

Сохраняем обычный первый порядок и хиральную нечётность. Вершина задаётся
тройкой
\begin{equation}
 v=(i,j,\chi),\qquad
 i,j\in\{\mathbb C,\mathbb H,M_3\},\qquad
 \chi\in\{L,R\}.
 \label{eq:v7-r2-branch-typed-vertex}
\end{equation}
Строго допустимое ребро удовлетворяет одновременно
\begin{equation}
 \chi_v\ne\chi_w,\qquad
 i_v=i_w\quad\text{или}\quad j_v=j_w.
 \label{eq:v7-r2-branch-strict-edge}
\end{equation}

Машинный перебор всех $18$ типов показывает: без новых вершин и с одной
новой вершиной цикла, содержащего оба существующих ребра
$Q_L-u_R$ и $L_L-e_R$, нет. Минимум равен
\begin{equation}
 N_{\rm new}^{\rm strict}=2.
 \label{eq:v7-r2-branch-minimum-new-vertices}
\end{equation}
Существует девять минимальных пар координат. Наименьшая по сырой комплексной
размерности добавляет зеркальные лептонные типы
\begin{equation}
 X_L=(\mathbb C,\mathbb C,L),\qquad
 Y_R=(\mathbb H,\mathbb C,R),\qquad
 \dim_{\mathbb C}(X_L\oplus Y_R)=1+2=3,
 \label{eq:v7-r2-branch-minimal-mirror-pair}
\end{equation}
и даёт шестикромочный цикл
\begin{equation}
 Q_L\to u_R\to X_L\to e_R\to L_L\to Y_R\to Q_L.
 \label{eq:v7-r2-branch-six-cycle}
\end{equation}
Здесь смешанный эффект факторизован по нескольким допустимым полям; прямого
элементарного ребра $R_2$ всё равно нет. Кроме того, ещё не выполнены
Real-удвоение, аномальный аудит, спектральное действие и экспериментальные
границы на зеркальных фермионов.

\section{Решение развилки}

Две оставшиеся ветви имеют несравнимую цену. Обобщённая флуктуация сохраняет
фермионный набор, но отказывается от линейного первого порядка и вводит
$A_{(2)}$. Строгий маршрут сохраняет аксиому, но требует как минимум двух
новых хиральных вершин и не создаёт элементарный $R_2$. Поэтому без заранее
заданного критерия стоимости нельзя честно объявить одну ветвь канонической.

\begin{theorem}[Минимальная архитектурная цена]
На неизменённой $\mathcal A_{\rm SM}$ автоморфная скрутка не переставляет
слагаемые и не лечит диагональный дефект $R_2$. Маршрут без новых фермионов
требует отказа от строгого первого порядка и квадратичного $A_{(2)}$.
Маршрут со строгим первым порядком требует не менее двух новых хиральных
вершин; существует ровно девять минимальных координатных пар.
\end{theorem}

\begin{nogocriterion}[Запрет ложного выбора]
Нельзя выбирать обобщённую или зеркальную ветвь только по меньшему числу
символов. Первая должна доказать конкретную опору $A_{(2)}$, Real-замыкание
и цветовой вакуум; вторая --- Real/аномальное завершение и отсутствие
произвольных новых юкавских матриц.
\end{nogocriterion}

\begin{protocol}
\begin{enumerate}
 \item простая автоморфная скрутка $\mathcal A_{\rm SM}$: \textbf{закрыта};
 \item обобщённая флуктуация, новые фермионные вершины: \textbf{ноль};
 \item необходимый новый член: \textbf{$A_{(2)}$};
 \item строгий путь, минимум новых вершин: \textbf{две};
 \item минимальных координатных пар: \textbf{девять};
 \item минимальная сырая комплексная размерность: \textbf{три};
 \item элементарный строгий $R_2$: \textbf{не получен};
 \item единственная физическая ветвь выбрана: \textbf{нет}.
\end{enumerate}
\end{protocol}

Машинный сертификат сохранён в
\path{s2t_v7_r2_minimal_architecture_branch_gate_results.json}.

\begin{thebibliography}{9}
\bibitem{v7r2branch-generalized}
A.~H.~Chamseddine, A.~Connes, W.~D.~van Suijlekom,
\emph{Inner Fluctuations in Noncommutative Geometry without the First Order Condition},
Journal of Geometry and Physics 73 (2013), 222--234; arXiv:1304.7583.
\bibitem{v7r2branch-landi}
G.~Landi, P.~Martinetti,
\emph{On twisting real spectral triples by algebra automorphisms},
Letters in Mathematical Physics 106 (2016), 1499--1530; arXiv:1601.00219.
\bibitem{v7r2branch-grand}
A.~Devastato, P.~Martinetti,
\emph{Twisted spectral triple for the Standard Model and spontaneous breaking of the Grand Symmetry},
Mathematical Physics, Analysis and Geometry 20 (2017), article 2;
arXiv:1411.1320.
\end{thebibliography}